% ===== 原文 pp.35–36 =====
% =====================================================================
%  tr05b —— 原文 pp.35–36（§4.2 后半：4.2.3 末 + 4.2.4）
%  编号覆盖：猜想 4.14、定理 4.15、定理 4.16、定理 4.17、定理 4.18
%  起点：p35 开头（接 tr05a 末段）；终点：p36 末（定理 4.18），p37 由 tr05c.tex 接
% =====================================================================

% [原文 p.35]
由上述内容可以看出，$\SO(p+q,p)$-Higgs 丛的 $(\mathrm{B},\mathrm{A},\mathrm{A})$-膜与 $\SO(2p+q,\C)$ 的 Hitchin 纤维化在点 $a$ 处的纤维之交，等同于两个模空间之积的一个覆盖：
\[
    \cM_{\mathrm{Cay}}(a) \times \cM_{\mathrm{Lan}}(a).
\]
该覆盖对应于上面 (iii) 中的扩张数据 $\tau$。Cayley 模空间 $\cM_{\mathrm{Cay}}(a)$ 与 Langlands 模空间 $\cM_{\mathrm{Lan}}(a)$ 如下给出：
\begin{itemize}[leftmargin=1.4em,itemsep=1pt,topsep=1pt]
\item $\cM_{\mathrm{Cay}}(a)$ 是 $\Sigma$ 上 $K^2$-扭（$K^2$-twisted）$\Gl(p,\R)$-Higgs 丛的模空间的 Hitchin 映射的一条纤维，可与上面 (i) 中的 $\Jac(\bar S)[2]$ 等同。

\item $\cM_{\mathrm{Lan}}(a)$ 是上面 (ii) 中 $\cC$ 上等变 $\SO(q)$-丛 $M$ 的模空间。
\end{itemize}

正如 [11] 所解释的，从谱数据重构 $\SO(p+q,p)$-Higgs 丛 $(E,\Phi)$ 需要取如下形式的扩张
\[
    0 \to V_0 \to E \to F\otimes K^{1/2} \to 0,
\]
其中 $F$ 是与 Cayley 模空间相联系的 $\Sp(2p,\R)$-Higgs 丛，$V_0$ 是 $\cC$ 上等变正交丛 $M$ 在 $\pi_{\cC} : \cC \to \Sigma$ 下的不变直像。扩张数据 $\tau$ 由定义上述扩张的扩张类给出，由此得到如下猜想。

\begin{conj}{4.14} Guichard 与 Wienhard 在 [53, Conjecture 5.6] 中猜想：$\SO(p+q,q)$-Higgs 丛模空间中那些含有正表示（positive representation）的额外连通分支，正是那些包含谱数据 $(L,M,\tau)$ 形如 $(\bar\pi^*I, \cO^{\oplus q}, \tau)$ 的 Higgs 丛的连通分支，其中 $I$ 是黎曼面 $\Sigma$ 上的非平凡 $\Z_2$-线丛。 \end{conj}
% [Tutor Note: p.35 猜想 4.14 作 $\SO(p+q,q)$-Higgs 丛，而本小节前文（p.34–p.35）一律为
%  $\SO(p+q,p)$；两种群不同，此处疑为笔误但无法仅据本讲否定原文，故照原文排，请对照 [53, Conjecture 5.6]。]

\para{4.2.4.\ 来自黎曼面上实结构的 $(\mathrm{A},\mathrm{B},\mathrm{A})$-膜} 如前所述，给定实结构 $f : \Sigma \to \Sigma$，存在一个由 [15, 17] 描述的自然的 $(\mathrm{A},\mathrm{B},\mathrm{A})$-膜，它作为诱导对合
\[
    i_2(\bar\partial_A,\Phi) = (f^*(\partial_A), f^*(\Phi^*)) = (f^*(\rho(\bar\partial_A)), -f^*(\rho(\Phi)))
\]
的不动点集出现。
% [Tutor Note: p.35 式中的第二个分量写作 $f^*(\partial_A)$，但同一等式的后一式与首项均作
%  $\bar\partial_A$（见 p.30 的式 (4.6)，此处与 (4.6) 完全相同），故 $f^*(\partial_A)$ 应作
%  $f^*(\bar\partial_A)$（漏印横杠）。照原文排，请读者注意。]

这些膜可以通过 [61] 引入的、与 $G_{\C}$-Higgs 丛相联系的谱数据来研究，并把 $(\mathrm{A},\mathrm{B},\mathrm{A})$-膜视为位于相应的 Hitchin 纤维化之中。这样做，[15] 便在 Hitchin 纤维化底的一个实子空间上得到了实可积系统的新例子：

\begin{thm}{4.15（[15] 定理 17）} 若 $(\mathrm{A},\mathrm{B},\mathrm{A})$-膜含有光滑点，则该膜与 Hitchin 纤维化之交给出一个带奇点的拉格朗日纤维化。其一般纤维是光滑的，由有限多个环面组成。 \end{thm}

当 $G_{\C} = GL(n,\C)$ 时，定理 4.15 中每条纤维的分支数可以用谱曲线上某个诱导对合的不动点集的分支数来表示（[15, Proposition 32]）。此外，在秩 2 的情形，计算可归结为用相应的二次微分（quadratic differential）给出的精确公式：

% [原文 p.36]
\begin{thm}{4.16（[15]）} 设 $q \in H^0(\Sigma,K^2)$ 是只有单零点的二次微分，且满足 $f^*(\bar q) = q$，并考虑谱曲线
\[
    S = \{0 = \eta^2 - q\} \subset \operatorname{Tot}(K).
\]
设 $n_+$ 为 $f$ 的、在其上 $q$ 不为零且为正的不动分支数，$u$ 为被 $f$ 固定的 $q$ 的零点数。则 $\Gl(2,\C)$ Hitchin 纤维化的一般纤维 $\Jac(S)$ 上诱导对合的分支数——也就是给出 $(\mathrm{A},\mathrm{B},\mathrm{A})$-膜的分支数——为 $2^d$，其中
\[
    d = \left\{\begin{array}{ll}
        2n_+ + \dfrac{u}{2} - 1 & \text{若 } 2n_+ + \dfrac{u}{2} > 0,\\[2ex]
        1 & \text{其他情形}.
    \end{array}\right.
\]
\end{thm}

此外，对于 $\Sl(2,\C)$，当考虑 Prym 簇 $\Prym(S,\Sigma)$ 时还可进一步精确化，见 [15, Theorem 38]。
% [Tutor Note: p.36 此句原文 "Moreover, for SL(2,C) can be specified further…" 缺少主语
%  （应为 "the statement can be specified further" 之类），属漏字；译文按正意补出"还可进一步精确化"。]

\begin{thm}{4.17（[15]）} 在定理 4.16 的设定下，并假设 $S$ 至少有一个分支点被 $f$ 固定。则 $\tilde f$ 在 $\Prym(S,\Sigma)$ 上作用的不动点集有
\[
    2^{n_0+\frac{u}{2}-1}
\]
个连通分支，其中 $n_0$ 是不含分支点的 $f : \Sigma \to \Sigma$ 的不动分支数，$u$ 是被 $f$ 固定的分支点数。 \end{thm}

研究这些 $(\mathrm{A},\mathrm{B},\mathrm{A})$-膜使我们能够考虑如下问题：给定一个以 $\Sigma$ 为边的 3-流形 $M$，把表示 $\pi_1(M)$ 限制到边界上便给出一个自然映射
\[
    \Rep(\pi_1(M),G_{\C}) \to \Rep(\pi_1(\Sigma),G_{\C}),
\]
并且可以证明，能够从边界延拓到该 3-流形上的表示总在 $\cM_{G_{\C}}$ 内定义一个 $(\mathrm{A},\mathrm{B},\mathrm{A})$-膜。虽然对这类由 3-流形产生的 $(\mathrm{A},\mathrm{B},\mathrm{A})$-膜尚无一般性描述，但理解它们的第一步是考察 [17] 中作为经由 (4.6) 的三元组的一部分而构造、并被 (4.6) 中的对合 $i_2$ 固定的那些 $(\mathrm{A},\mathrm{B},\mathrm{A})$-膜。事实上，考虑 3-流形
\[
    M = \Sigma\times[-1,1]/(x,t)\mapsto(f(x),-t),
\]
（任何柄体（handle body）都可写成这种形式），可以看出商 $M$ 是一个以 $\partial M = \Sigma$ 为边的 3-流形。
% [Tutor Note: p.36 "consider the the 3-manifolds" 中 "the" 重复；同段 "appearing though 3-manifolds"
%  中 "though" 应为 "through"。两处均为排印笔误，译文按正意译，原文照录于此备查。]

\begin{thm}{4.18（[17]）} 设 $(E,\Phi)$ 是 $i_2$ 的、具有单和乐（simple holonomy）的不动点。则相应的联络作为平坦联络延拓到 $M$ 上，当且仅当约化等变 K-理论中的类 $[E] \in \widetilde{K}^0_{\Z_2}(\Sigma)$ 是平凡的。 \end{thm}
